\section{Two further Tate residues for the original conductor-four character}
\label{sec:conductor4-quartic-residue-strict}

\paragraph{Inputs and scope.}
This proof uses the class-sector conclusion of AC-G152-class and
its exact field and unit data.  The class certificate, this strict
argument, and the finite local computations have separate recorded
review scopes.  Every rational denominator is retained below.
No original finite host character is enlarged: the conductor-$269$
quartic character appears only in the residual Kummer splitting field.

Put $p=269$, $K=\mathbb Q(\omega)$ and $t=\sqrt[3]p>0$.
Let $\chi:G_K\to\mu_3$ describe $K(t)/K$, and write
$\psi_4=\chi_{-4}|_{G_K}$.
Let $c_p$ be the actual cyclotomic character and $\omega_p$ its
Teichmuller part.  Fix $i\in\mathbb Z_p$ with $i^2=-1$ and
$\bar i=82\in\mathbb F_{269}$.
For a discrete coefficient module $D$ use the full strict condition
\[
 \begin{gathered}
 S_{\rm str}(K,D)=\ker\bigl(H^1_{\rm cts}(G_K,D)
                                      \longrightarrow\mathcal L(D)\bigr),\\
 \mathcal L(D)=\prod_{v\mid p}H^1(G_{K_v},D)
             \times\prod_{v\nmid p,\infty}H^1(I_v,D).
 \end{gathered}
 \tag{C4Q.1}
\]

\begin{theorem}[AC-G152: the conductor-four quarter-residue strict rows]
\label{AC-G152}
Let $E/\mathbb Q_p$ be finite and let $O$ be its original ring of
integers, containing the values of $\chi$.
Let $j=1,2$, $m\in\mathbb Z$, $e\geq1$, and let
$\kappa:G_K\to1+\mathfrak m_O$ be continuous and unramified
at every finite place away from $p$.
Set
\[
 D=(E/O)(\chi^{-j}\psi_4\omega_p^{-m}\kappa^{-1}c_p^{-e}).
 \tag{C4Q.2}
\]
If
$e+m\equiv66$ or $200\pmod{268}$, then
\[
                         S_{\rm str}(K,D)=0.
 \tag{C4Q.3}
\]
The corresponding residual strict groups vanish over every finite
field containing their character values.
No condition on the absolute ramification index of $O$ is imposed.
For $m=0,\kappa=1$, these are two new Tate progressions for the
already existing finite labels $\chi^{-j}\psi_4$.
\end{theorem}

\begin{proof}
\proofstep{1. The one real quartic Kummer field and both orientations.}
Let $\psi_p=\omega_p^{67}$ and
$\xi=\chi_{-4}\psi_p$ as rational Dirichlet--Galois characters.
Both factors are odd, so $\xi$ is even of order four.
Its conductor is $4p$, its square is
$\tau_p=\omega_p^{134}$, and $\psi_p(2)=-i$.
The exact character/period calculation in AC-G152-class identifies
its real cyclic quartic field with
\[
 \begin{gathered}
 F=\mathbb Q(w),\qquad w^4-pw^2+25p=0,\\
 L=F(t),\qquad M=KL.
 \end{gathered}
 \tag{C4Q.4}
\]
In particular
$w^2=(269\pm13\sqrt{269})/2$; both values are positive.
Use the generator
\[
               b(w)=\frac{w^3-219w}{65},\qquad \xi(b)=-i.
 \tag{C4Q.5}
\]
The fixed period orientation is verified by the CRT representative
$809$, which fixes $i$ and sends a primitive $p$th root of unity
to its square.  Its inverse has $\xi$-value $+i$.
The local argument below checks both quartic eigenvalues, rather
than relying on an interchange of these two orientations.

The unique quadratic subfield of $K(t)/\mathbb Q$ is the imaginary
field $K$, which cannot lie in the totally real field $F$.
As the possible common degree divides both four and six, their
intersection is $\mathbb Q$.
Consequently
\[
 \begin{gathered}
 \operatorname{Gal}(M/\mathbb Q)=S_3\times C_4,\qquad
 \Delta=\operatorname{Gal}(M/K)=C_3\times C_4,\\
 [L:\mathbb Q]=12,\qquad[M:\mathbb Q]=24.
 \end{gathered}
 \tag{C4Q.6}
\]
Let $c(t)=\omega t$ generate $C_3$, with $c(w)=w$;
let $b$ fix $K$ and $t$.  These two automorphisms commute.

Write $k=O/\mathfrak m_O$ and let $a$ be the residual character
in (C4Q.2).  Put $r\equiv e+m\pmod{268}$.
Its Kummer quotient is $\eta=a/\bar c_p$, and the precise table is
\[
 \begin{array}{c|c|c}
 r&\eta&\eta^{-1}(b)\\ \hline
 66&\bar\chi^{-j}\bar\xi^{-1}&-\bar i=187\\
 200&\bar\chi^{-j}\bar\xi&\bar i=82.
 \end{array}
 \tag{C4Q.7}
\]
Indeed $-(66+1)=-67$ and $-(200+1)\equiv67\pmod{268}$.
The character $\eta$ is faithful on $C_3\times C_4$.
Thus the same field $M$ splits it in both rows.

\proofstep{2. Exact Kummer descent and the needed class sector.}
The quotient order twelve is prime to $p$.
Restriction and inflation, including the local decomposition groups,
therefore identify the residual strict group with
\[
 \begin{gathered}
 e_{\eta^{-1}}\ker\left(V_{M,S}\otimes_{\mathbb F_p}k
       \longrightarrow\mathcal P_M\otimes_{\mathbb F_p}k\right),\\
 \mathcal P_M=\bigoplus_{v\mid p}M_v^\times/M_v^{\times p}.
 \end{gathered}
 \tag{C4Q.8}
\]
where
$V_{M,S}=\{x\in M^\times:p\mid v(x)\text{ for }v\nmid p\}/M^{\times p}$.
Over $M$ the coefficient becomes $k(1)$.
The additional quotient action on its Kummer cohomology is the twist
$\eta$, so invariants select $\eta^{-1}$.
The full local power condition also forces the valuations at the
primes above $p$ to be divisible by $p$.
Thus the kernel lies in the all-primes valuation group $V_M$, with
\[
 0\longrightarrow\mathcal O_M^\times/\mathcal O_M^{\times p}
 \longrightarrow V_M\longrightarrow\operatorname{Cl}(M)[p]
 \longrightarrow0.
 \tag{C4Q.9}
\]

The class input is precisely $\operatorname{Cl}(L)[p]=0$,
proved by AC-G152-class.  Let $s$ be the reflection in the
$S_3$ factor fixing $L$.
Norm to $L$ followed by extension acts on $\operatorname{Cl}(M)[p]$
as $1+s$, and is zero.  Hence $s=-1$ on this module.
The same holds for every conjugate reflection.  Products of these
reflections generate $C_3$, so $C_3$ acts trivially.
Only its nontrivial cubic sector is needed in (C4Q.8), and that
sector vanishes.  No assertion about the whole class group of $M$
is made or used.

\proofstep{3. The unique completion and an exact uniformizer.}
The quartic polynomial in (C4Q.4) is Eisenstein at $p$.
The quartic and cubic completions are totally ramified of coprime
degrees four and three.  Their compositum has degree and
ramification degree twelve.
The completion $K_p/\mathbb Q_p$ is unramified quadratic.
Adjoining it gives the full global degree twenty-four of $M$.
Consequently $M$ has exactly one prime above $p$, of absolute
ramification index twelve and residue degree two.

In that completion take $u=t/w$, so that $v_u(t)=4$,
$v_u(w)=3$ and $v_u(u)=1$.
Elimination, equivalently the following exact inverse formulas, gives
\[
 \begin{aligned}
 15625u^{12}-52186u^6+269&=0,\\
 w&=\frac{52811u^3-15625u^9}{244},\qquad t=uw.
 \end{aligned}
 \tag{C4Q.10}
\]
Direct substitution and reduction by the first line verify
$t^3=p$ and $w^4-pw^2+25p=0$.
The polynomial after division by $15625$ is Eisenstein at $p$,
and the inverse formulas identify the resulting degree-twelve
field with the completion of $L$.
Its power basis is the full integral basis locally.
Indeed the valuations of the twelve terms $a_nu^n$ with
$a_n\in\mathbb Q_p$ are distinct modulo twelve, so the expansion
is integral exactly when all $a_n$ are $p$-integral.
The same argument works over the unramified coefficient ring
$\mathcal O_{K_p}$ for the completion of $M$.

The exact action induced by (C4Q.5) is
\[
 \begin{aligned}
 c(u)&=\omega u,\\
 b(u)&=\frac{46711}{15860}u-\frac{3125}{3172}u^7
       \equiv82u+51u^7\pmod{p,u^{12}}.
 \end{aligned}
 \tag{C4Q.11}
\]
Substitution verifies $b(w)$ as in (C4Q.5), $b(t)=t$,
and $b^2(u)=-u$ exactly.
Every denominator in (C4Q.10)--(C4Q.11) is prime to $p$.
These are actual actions on the unique completion, not guessed
linear actions on a chosen quotient.

\proofstep{4. An actual integral unit, with all denominators retained.}
Use unit number six in the exact seven-unit file of AC-G152-class.
In the ordered tensor basis $t^aw^b$, $0\leq a<3$, $0\leq b<4$,
with coordinate index $4a+b$, it is the vector below divided by
$52455=195p$:
\[
 \begin{aligned}
 (&54808414762075887245245,
       3530543866123251071586,\\
  &-1965017582145838334080,
       -126578752579972745404,\\
  &8492433594955935699440,
       546791606561600990739,\\
  &-304474803763916613125,
       -19603835019267409661,\\
  &1315407208379945993845,
       84714730690320710955,\\
  &-47160610343660095945,
       -3037233169320092320).
 \end{aligned}
 \tag{C4Q.12}
\]
The complete monic integral degree-twelve relation with constant
coefficient one is included in
\path{selected_conductor4_unit.json}, whose SHA256 is
\begin{quote}\ttfamily\small
839f9999bb007912b5b221564f261b236\\
332b9235e0775fe039be6af989b728a.
\end{quote}
The independent verifier substitutes (C4Q.10) into the entire
rational tensor expression and checks that relation by Horner
evaluation in $\mathbb Q[u]$ modulo (C4Q.10).
It checks all seven supplied unit relations, although only (C4Q.12)
is required here.  Monicity proves integrality, and the constant
coefficient one expresses the inverse as an integral polynomial
in the unit.  No complete fundamental-unit basis is needed.

Only after the rational change of basis does the verifier reduce
the coefficients modulo $p$.
In particular it retains the factor $p$ in $195p$ until cancellation.
All twelve resulting rational coefficients are $p$-integral.
Their exact reduction, in the order $1,u,\ldots,u^{11}$, is
\[
 \begin{aligned}
 U\equiv(&1,21,86,47,248,155,\\
         &113,81,258,186,45,33)\pmod{p,u^{12}}.
 \end{aligned}
 \tag{C4Q.13}
\]
Thus $U$ is a principal unit at the unique prime above $p$.

\proofstep{5. Both quartic projectors are nonzero.}
The complete logarithm of (C4Q.13), truncated below degree twelve, is
\[
 \log U=21u+252u^3+152u^5+128u^7+131u^9+37u^{11}.
 \tag{C4Q.14}
\]
All logarithmic denominators one through eleven are invertible
modulo $p$.  The verifier computes this logarithm independently
by solving $U(\log U)'=U'$ recursively and integrating termwise.
Form the actual global cubic norm correction
$V=U^3/(U\,c(U)\,c^2(U))$.
The operator $3-(1+c+c^2)$ kills precisely the exponents divisible
by three and multiplies the others by three.  It gives
\[
              \log V=63u+187u^5+115u^7+111u^{11}.
 \tag{C4Q.15}
\]
For $\lambda=82,187$, let
$P_\lambda=\sum_{r=0}^3\lambda^{-r}b^r$ on the unit classes
modulo $p$.
This is four times the corresponding idempotent; no division
by $p$ occurs.  Applying the full action (C4Q.11) gives
\[
 \begin{aligned}
 P_{82}\log V&=252u+210u^5+39u^7+28u^{11},\\
 P_{187}\log V&=152u^7+147u^{11}.
 \end{aligned}
 \tag{C4Q.16}
\]
Both are nonzero.  The exact computation also verifies
$b(P_\lambda\log V)=\lambda P_\lambda\log V$.
The coefficients of $P_\lambda$ may be represented by integers;
the corresponding product of the global units $b^r(V)$ realizes
the projected Kummer class before localization.
Changing a representing integer by a multiple of $p$ does not
change that class or the displayed truncated logarithm.

For the local field at hand, $v_u(p)=12<p-1$.
Logarithm and exponential are inverse on its principal units:
the nonlinear terms have strictly larger valuation.
The logarithm of a $p$th power principal unit lies in the thirteenth
power of the maximal ideal.  The nonzero terms of degrees one
and seven in (C4Q.16) therefore detect two nonzero full local
Kummer classes.  If a principal unit is a $p$th power of a unit,
its root also has residue one, since Frobenius is bijective on
the finite residue field.  Thus no larger local power subgroup
has been overlooked.

\proofstep{6. Unit multiplicities and residual strict zero.}
The twelve archimedean places over a fixed complex embedding of
$K$ form a regular $\Delta$-set.
The equivariant Dirichlet unit character is
$\mathbb Q[\Delta]-\mathbb Q$.
There are no roots of unity of order $p$ in $M$, since
$268$ does not divide $[M:\mathbb Q]=24$.
Since $p\nmid12$, reduction preserves these character multiplicities.
For each $b$-eigenvalue $\lambda=82,187$, the nontrivial cubic
part of the global unit quotient has dimension two over
$\mathbb F_p$.
It is simple as an $\mathbb F_p[C_3]$-module, because
$X^2+X+1$ is irreducible over $\mathbb F_{269}$.

Localization is equivariant for the entire group $\Delta$.
The two nonzero values in (C4Q.16) show that its restriction to
each of these simple two-dimensional spaces is injective.
After extension to $k$, both cubic eigenlines in both rows of
(C4Q.7) remain injective.
The vanished nontrivial cubic class sector and (C4Q.9) now prove
that the kernel in (C4Q.8) is zero.
This proves the residual assertion, with the full local condition.

\proofstep{7. Lift on the original coefficient ring.}
Let $\varpi$ be a uniformizer of the original $O$.
The module $D$ is divisible and $D[\varpi]\simeq k(a)$.
At $p$, $\psi_4$ is unramified, $\kappa$ reduces to one,
and the residual cubic inertia has order three.
A residual cyclotomic character has order dividing $268$,
so cannot cancel that cubic character.
Some inertia element therefore acts by a scalar whose difference
from one is a unit of $O$; hence $D^{G_{K_p}}=0$.
At every finite place away from $p$, the cyclotomic, Teichmuller
and $\kappa$ factors are unramified.
The remaining finite inertia character has order dividing six,
prime to $p$, so its invariants are zero or the whole divisible
module $D$.
For every local group $J$ in (C4Q.1),
$D^J/\varpi D^J=0$.

The multiplication-by-$\varpi$ long exact sequence gives local
injectivity $H^1(J,D[\varpi])\to H^1(J,D)$.
A global strict class killed by $\varpi$ lifts to a residual
class, and local injectivity makes that lift strict.
It vanishes by Step6.
Every continuous cocycle from compact $G_K$ to the discrete
module $E/O$ has finite image and is killed by a power of
$\varpi$.  A nonzero strict class would therefore yield a
nonzero one killed by $\varpi$, a contradiction.
This proves (C4Q.3) without replacing $O$ or assuming it is
unramified over $\mathbb Z_p$.
\end{proof}

\paragraph{Actual host labels and the exact new weight count.}
For $m=0,\kappa=1$ the finite character is still
$\chi^{-j}\chi_{-4}$, which occurs on every actual AC-K25 host
containing $i$ and $K(t)$.
Neither $\psi_p$ nor $\tau_p$ is added to an original finite
host label; they only describe the residual splitting field.
For positive integer $T$, the new weights are $e=66+134n$,
$n\geq0$, and have total count
$2\lfloor(T+67)/134\rfloor$ for the two cubic labels.
Together with the already proved weights $133,267\pmod{268}$,
the same two finite labels have the full progression
$e\equiv-1\pmod{67}$ and count $2\lfloor T/67\rfloor$.
This does not itself set any positive derived index to zero.
The full local comparison cost and common-window transport
require their separate argument.

\paragraph{Evidence and geometric scope.}
The complete new finite calculation is
\path{verify_conductor4_quartic_local.py} and its JSON output.
It retains the full source unit arrays, all rational changes of
basis, the monic relations, and both projected local values.
AC-G152-class is used only for its precise field/character
dictionary, actual unit arrays and $\operatorname{Cl}(L)[p]=0$.
No full class group of $M$, complete unit basis, or analytic
interval rerun is asserted by the local verifier.
Kummer theory uses the previously audited Stacks tag~03PK input.
Finite prime-to-$p$ descent, equivariant Dirichlet units and
class extension/norm retain their earlier checked scopes.
SGA~1, Expos\'e~V, \S\S4 and 7 supplies the retained finite-cover
and rank-one Artin--Tate dictionary only.
No automorphic endpoint, mixed realization, all-prime estimate,
H04.4 conclusion or Lean execution is supplied by this proof.
